\paragraph{\fid{F55}, Part 1: a variant retaining the actual formal bar.}
The one-place construction can additionally commute with
an explicit extension of the actual tensor bar.
This variant retains the actual right action and both
trace laws, while still failing all-places regularity
and the stronger left/word comparisons.

\begin{proof}
Keep the actual bimodule decomposition and regular
bar-fixed basis of \fid{F43}, and the constants
$s=\sqrt2$, $\eta=189/625$, $v=t^2$ from Part~1.
Extend the coefficient involution by
\[
 \overline t=t^{-1},\qquad \overline s=s.
\]
This is a stated choice of lift of $v\mapsto v^{-1}$
to the coefficient cover.
Define a second family, distinguished by tildes:
\[
 \widetilde g(t)=\frac{2(1+t)^2}{t}
                  =4+2t+2t^{-1},\qquad
 \widetilde h(t)=\frac{\eta\widetilde g(t)}{\Delta(t^2)},
\]
\begin{equation}\label{eq:barfractionalfunctions}
 \widetilde\lambda(t)=\frac{\widetilde h(t)}2-2-s,\qquad
 \widetilde u(t)=1+\frac{\widetilde g(t)-\widetilde h(t)}{20}.
\end{equation}
All four functions are fixed by this involution:
$\widetilde g(t^{-1})=\widetilde g(t)$ and
$\Delta(v^{-1})=\Delta(v)$.
The cancellation formula is
\[
 \widetilde h(t)=
 \frac{2\eta(t^4+1)^4}
 {t(t-1)^2(t^2+1)^2(t^8+t^4+1)}.
\]
The exact special values are
\[
\begin{array}{c|cccc}
t&\widetilde g&\widetilde h&\widetilde u&
                                   \widetilde\lambda\\ \hline
s&4+3s&4+3s&1&s^{-1}\\
-1&0&-126/625&6313/6250&-1313/625-s .
\end{array}
\]

Use the same $R_0$ as in Part~1 and set
\[
 \widetilde R
   =R_0[\widetilde u^{-1},\widetilde\lambda^{-1}].
\]
It is stable under bar: the inherited factors are
bar invariant, $t-1$ maps to $-(t-1)/t$,
$t^2+1$ to $(t^2+1)/t^2$, and the two new functions
are invariant.
Every inherited denominator is nonzero at $t=-1$
and all $t_n=s^n$.
The displayed values check the additional units
at $-1$ and $s$.
For $n>1$, $\Delta(2^n)>\eta$, so
\[
 \widetilde h(t_n)<\widetilde g(t_n),\qquad
 \widetilde u(t_n)>1.
\]
If $n$ is even, $t_n$ is rational and
$\widetilde\lambda(t_n)\in\Q-s$, hence is nonzero.
For $n=2r+1$ odd, put
\[
 d_n=\Delta(2^n),\qquad c_r=2^r+2^{-r-1}.
\]
Both are rational, and
\[
 \widetilde\lambda(t_n)
 =2\left(\frac{\eta}{d_n}-1\right)
   +s\left(\frac{\eta c_r}{d_n}-1\right).
\]
At $n=1$ this equals $s^{-1}$.
For odd $n>1$, its rational coefficient of $1$
is nonzero because $d_n>\eta$.
Linear independence of $1,s$ proves nonvanishing.
Thus one bar-stable finite localization contains the
chosen collision and every extension point.
The contact order at $t=-1$ is still two.

On the entire rank-$20$ $\epsilon_+$ block use
$\widetilde u I_{20}$, and on the rank-$36$
$\epsilon_-^4\oplus N(\delta)^{16}$ block use identity.
On the final multiplicity factor, in order
$(++,+-,-+,--)$, put
\[
 \widetilde F
   =\operatorname{diag}(1,1,s,\widetilde\lambda)\otimes I_2.
\]
Its determinant is
$2\widetilde u^{20}\widetilde\lambda^2\in\widetilde R^\times$.
As before, transport by
$\iota^{-1}\widetilde F_{\rm block}\iota$ to the
actual tensor lattice.
The operator commutes with the actual right action
because it acts on multiplicity factors.
Its entries and all adapted basis vectors are fixed
by the extended actual bar $\beta$, so
\[
 \beta\widetilde F=\widetilde F\beta
\]
as semilinear maps over $\widetilde R$.
The established bar-compatible basis comparison
retains this equality on the original tensor module.

The marked trace is
$2(1+1+s+\widetilde\lambda)=\widetilde h$.
Hence
\begin{equation}\label{eq:barfractionaltraces}
 \Tr(X(t^2)\widetilde F(t))=\eta\widetilde g(t),\qquad
 \Tr\widetilde F(t)=56+\widetilde g(t).
\end{equation}
At $t=s$, the marked eigenvalues $1,s,s^{-1}$
have multiplicities $4,2,2$; on the full module
the multiplicity of $1$ is $60$.
Thus
\[
 \Tr(\widetilde F(s)^n)
       =60+2s^n+2s^{-n}=56+\widetilde g(t_n),
\]
\[
 \Tr(X(2)\widetilde F(s)^n)
       =\eta(4+2s^n+2s^{-n})=\eta\widetilde g(t_n).
\]
These give both retuning identities for every $n$.

The stronger comparisons still fail explicitly.
On the final multiplicity space,
\[
 (\widetilde F(s)E_V-E_V\widetilde F(s))
       (g_-\otimes g_+)
  =(1-s)(g_+\otimes g_+)\ne0.
\]
The generator calculation gives
\[
 \Tr(P_1(v)\widetilde F(t))
    =D(v)\left(36+\widetilde g(t)
                         -\frac{\widetilde h(t)}2\right),
\]
\[
 \Tr(P_1(2)\widetilde F(s)^n)
    =D(2)\left(38+s^n+s^{-n}\right)
    =D(2)\left(36+\frac{\widetilde g(t_n)}2\right).
\]
At $n>1$ the first is strictly larger, since
$\widetilde h(t_n)<\widetilde g(t_n)$ and $D(2^n)>D(2)$.
Multiplication by the respective $D$ factors proves
the same failure for $P_1^2$.

Unlabelled spectral retuning fails as well:
$\widetilde F(s)^n$ has eigenvalue $1$ of multiplicity
$60$, whereas $\widetilde F(t_n)$ has it with
multiplicity at most $42$ for $n>1$.
The latter count is $36$ unchanged dimensions,
four fixed marked dimensions, and at most two more
if $\widetilde\lambda(t_n)=1$; neither
$\widetilde u(t_n)>1$ nor $s>1$ contributes one.

Finally, $\widetilde h,\widetilde u,
\widetilde\lambda$ have order-two poles at $t=+1$,
with the latter two leading terms
$-\widetilde h/20$ and $\widetilde h/2$.
These are actual eigenvalue poles and cannot be
removed by a finite free lattice change.
The involution fixes $t=+1$ and $t=-1$ separately.
Bar compatibility consequently does not supply
all-places regularity or continuation from the
positive arithmetic sheet to the negative collision.
\end{proof}

At $q=4$ under $v=q^{1/2}$, this variant's marked
weights are $0,\tfrac12,-\tfrac12$.
It is not an integral mixed-geometric or pure-total
realization.
The formal bar covers $t\mapsto t^{-1}$; it is not
asserted to be a Frobenius-equivariant identification
of reciprocal weight lines at one fixed arithmetic
parameter, nor Verdier duality or VIII globalization.
The original Part~1 model and its comparisons are
unchanged.

% Separate paragraph for Part III.
\paragraph{\fid{F55}, the bar-compatible variant, in Grothendieck's language.}
The coefficient cover carries a specified involution
lifting inversion on the parameter torus, and the actual
coherent module automorphism is equivariant for this
involution and the actual tensor bar.
The affine module framework is EGA~I, Theorem~1.3.7,
p.~86, Corollary~1.3.8, p.~87, and Theorem~1.4.1,
pp.~90--92.
The representation/coaction framework is SGA~3, I,
4.7.2--3, with cote~90, scan~17 as the audited anchor
\cite{EGA,SGA3,Archive}.
The involution on the algebraic coefficient cover is
explicit additional equivariance data.
It does not select VIII's prescribed global integral
form or convert this mixed algebraic spectrum into
an integral-weight arithmetic realization.
The actual bar-fixed adapted basis is the proved
input from \fid{F43}; the invariant operator family
and all unit and trace identities are proved above.

